\section{Discussion and Conclusion}
\label{sec:conclusion}
\label{sec:discussion}

\aknew{\paragraph{IDS is language- and problem-agnostic.} IDS as a technique depends on neither the verification backend nor the problem domain. The backend need only support (i) mechanical checking of partial proofs, (ii) a programmatic interface an LLM agent can drive, and (iii) executable extraction of verified implementations for performance feedback; Lean~\cite{lean} and Verus~\cite{verus} both qualify, and we expect IDS to transfer to either without structural changes. The problem domain need only admit a machine-checkable formal specification and a way to benchmark candidate implementations. We chose \coq{} for the maturity of its ecosystem, and we chose distributed key-value stores because the literature already provided a complete formal specification for one consistency model; following its structure, we contribute six additional specifications covering other consistency properties (\cref{app:suite}).}

\aknew{\paragraph{Limitations. \ } Our work has three main limitations. First, \sys{} requires a formal \coq{} specification as input; writing complete formal specifications for distributed systems requires substantial expert effort, a challenge we return to below. Second, the seven specifications we evaluate cover correctness-critical behavior but not the full operational envelope of a production KV store; they do not define interfaces for scaling out (adding or removing nodes at runtime), reconfiguration, fault recovery, or informal but practically necessary behavior such as logging and observability endpoints. Given a richer specification that captures these aspects, we see no reason IDS would not apply unchanged, but evaluating this is left to future work. Third, our empirical evaluation covers the distributed key-value store problem only, leaving validation on other domains (e.g., OS protocols, cryptographic primitives) to future work. \cref{app:limitations-impacts} discusses additional limitations and broader impacts.}

\aknew{\paragraph{The specification bottleneck.} While \sys{} automates the synthesis of verified implementations, writing the formal specification itself remains the largest open problem: accurately capturing a system's intended behavior in a machine-checkable form is hard, and any correctness guarantee is only as strong as the specification it is checked against. We envision two directions for closing this gap. First, \emph{adversarial specification synthesis}, in which an LLM agent collaborates with a human stakeholder in natural language, iteratively probing intent and surfacing ambiguities until a formal specification emerges. Second, \emph{specification extraction from existing systems}, in which an agent recovers the correctness properties of a deployed implementation while leaving the design space open, so that \sys{} can synthesize alternative implementations that satisfy the same properties with better performance. We are actively pursuing both directions, which we plan to report on in future work.}

%Verified systems construction has long taken months to person-years of expert effort. \sys{} replaces that effort with compute, in hours: a proving-based regime where an LLM coding agent receives fast, exact correctness feedback with no false positives, turning joint code-and-proof construction into a search problem the agent can drive.

\paragraph{Conclusion. \ } Verified systems construction has traditionally required months to years of expert effort. \sys{} condenses this timeline to mere hours by replacing human labor with compute. By providing an LLM coding agent with fast and exact correctness feedback without false positives, \sys{} transforms the joint code-and-proof construction into an automated search problem. This shifts vibe coding to verified coding by using LLMs to directly build type-safe systems that are machine-checked against a formal specification, effectively eliminating silent bugs and unverified assumptions. Crucially, this methodology generalizes to any domain with a machine-checkable correctness oracle, including OS kernels, compilers, cryptographic protocols, and hardware. Ultimately, \sys{} shows that verified software generation is transitioning from a human-labor bottleneck to a compute-driven approach, paving the way for the broad adoption of verified system synthesis.

\section*{Acknowledgments}
This research is supported by NSF (IFML) CCF-2019844 and gifts from Accenture, AMD, Anyscale, Broadcom Inc., Google, IBM, Intel, Intesa Sanpaolo, Lambda, Mibura Inc., Samsung SDS, and SAP. We also thank our Sky Lab and NetSys colleagues at UC Berkeley for fruitful discussions that shaped this paper.

% proof-driven 
%Future extensions could also try to solve the spec-writing problem with agents and explore joint optimization across other resource axes (memory, energy, latency). Verified software is becoming a compute problem rather than a person-year problem, and adoption will only grow as LLMs are used for certified system building and proof-driven coding beyond what informal review can cover.
% Verified systems construction has long taken months to years of expert effort. IDS replaces that effort with compute in hours: by synthesizing both the implementation and its proof incrementally and in lockstep, IDS turns formal verification into a search problem. The agent receives fast feedback from a formal checker to rule out dead ends, while a distributed testbed evaluates throughput to steer the search toward efficient implementations. This enables agents to autonomously generate complex distributed systems with end-to-end formal correctness guarantees, avoiding the silent bugs that escape standard testing. The specification itself still has to be written by a person, but everything that follows can be discovered, and the recipe could apply wherever a machine-checkable correctness oracle exists (OS kernels, compilers, cryptographic protocols, hardware). Future extensions could try to solve the spec-writing problem with agents and explore joint optimization across other resource axes (memory, energy, latency). Verified software is becoming a compute problem rather than a manual problem, and adoption will only grow as LLMs are used for certified system building beyond what informal review can cover.